\section{Discussion}

The four models support a conservative synthesis.  Exactness regions are useful, but they do not have one authority or one explanation.  A theorem about the full feasible family can coexist with an incomplete named taxonomy.  A proof-derived search ceiling can be much larger than intrinsic support.  An objective region can describe where a proof works without describing what is necessary outside.  A feature vocabulary can determine a complete finite domain while carrying almost one cell per instance and failing at the next size.

The field-level consequence is a change in admissible comparison practice, not a universal performance prediction.  A support ceiling may be entered as a search restriction but cannot be compared as intrinsic expressivity until tightness is shown.  A proof-validity region may index a certificate but its complement cannot be labelled a necessity region.  Zero mixed cells may establish finite determination but cannot be reported as generalization without a transfer coordinate.  The present results supply a counterexample to each prohibited inference.  Recording the fields separately therefore changes which conclusions can be drawn from common exact-compilation outputs.

This separation clarifies the relation to algorithm selection.  A performance footprint answers which method wins on which features.  The present record asks additional questions when exact structure is available.  Is there a constructive witness for the gap?  Does every optimum admit a normal form?  Is the normal form tight?  Which objective weights preserve the proof?  Does the representation determine the label?  The answers can constrain or explain a footprint, but they do not replace instance-space analysis.

The state-preparation result gives the strongest warning against overcompression.  Eliminating mixed cells by adding sign-aware features might appear to solve the boundary problem.  The cell counts show otherwise.  Most complete-domain states become singletons, and transfer coverage collapses to two panel states.  The conversion is scientifically useful because it localizes the earlier failure to missing sign information.  It does not produce a reusable predictor.

The same logic applies to support.  The syndrome-rank ceiling remains useful for restricting search, even though it is not tight.  The tight support-one result required a richer proof move, namely global tag reconstruction.  Proof-language improvement, rather than larger enumeration alone, exposed the intrinsic value.

For practitioners, the reporting order is direct.  State the exact model and objective.  Record where the reference is exact.  Attach constructive witnesses to failures.  Separate intrinsic support from a safe proof ceiling.  Index every certificate by objective weights.  Test whether features determine the label before increasing classifier capacity.  Freeze transfer tests when an exact referee is available.  This order keeps favorable and adverse results on the same evidential footing.
